\subsection{استواری معماری سهم عروقی و تشخیص توجه}

دو پرسش این زیربخش جدا هستند. آیا سهم بازنمایی عروقی به یک استخراج‌کننده‌ی خاص ویژگی تصویری گره خورده است، و وزن‌های توجه این معماری‌ها تا چه اندازه قابل خواندن‌اند؟ هر دو پرسش روی آزمون تقسیم استاندارد قفل‌شده و با همان بوت‌استرپ زوجی طبقه‌بندی‌شده پاسخ داده می‌شوند.

برای پرسش نخست، همان نمایش منجمد عروق به دو بازنمایی تصویری متفاوت اضافه شد و هر سطر کنترل هم‌تای فقط‌تصویری خودش را دارد. سطر نخست بازنمایی اصلی تک‌ستون‌فقرات \lr{EfficientNet-B5} است و سطر دوم بازنمایی دو‌ستون‌فقرات توجهی که دو رمزگذار \lr{ResNet50} و \lr{EfficientNet-B4} را با توجه نرم ترکیب می‌کند. جدول زیر تفاضل‌ها را کنار هم می‌گذارد.

\begin{table}[htbp]
\centering
\renewcommand{\arraystretch}{1.35}
\footnotesize
\caption{استواری سهم عروقی نسبت به معماری بازنمایی تصویری}
\label{tab:archrobust}
\setlength{\tabcolsep}{3pt}
\resizebox{\textwidth}{!}{%
\begin{tabular}{|>{\raggedleft\arraybackslash}m{3.0cm}|r|r|r|r|r|r|r|}
\hline
بازنمایی تصویری & فقط تصویر & تصویر + عروق & تفاضل & کران پایین & کران بالا & \lr{p} & تفاضل \lr{Brier} \\
\hline
تک‌ستون‌فقرات اصلی (\lr{EfficientNet-B5}) & 0.9249 & 0.9328 & +0.007938 & $+0.002630$ & $+0.013131$ & 0.0031 & $-0.027693$ \\
\hline
دو‌ستون‌فقرات توجهی (\lr{ResNet50} + \lr{EfficientNet-B4}) & 0.9060 & 0.9269 & +0.020936 & $+0.011426$ & $+0.030828$ & $<0.0001$ & $-0.049181$ \\
\hline
\multicolumn{8}{p{0.96\textwidth}}{\footnotesize هر سطر یک آزمایش مستقل با کنترل هم‌تای خودش است و دو سطر یک مقایسه‌ی مدل واحد نیستند. هر دو سطر از همان نمایش تثبیت‌شده‌ی 1792بعدی عروق استفاده می‌کنند. اینکه سطر دوم افزایش بزرگ‌تری دارد به معنای برتری مدل دو‌ستون‌فقرات نیست؛ \lr{AUC} کنترل فقط‌تصویری آن در جدول پایین‌تر است.} \\
\hline
\end{tabular}}
\end{table}

افزودن عروق در هر دو سطر بازه‌ی بیرون از صفر و بهبود کیفیت احتمال می‌دهد. خواندن این الگو باید محدود بماند. بازتولید سهم عروقی با یک استخراج‌کننده‌ی دوم نشان می‌دهد این سهم به یک معماری تصویری خاص گره نخورده است، ولی به این معنا نیست که مدل دو‌ستون‌فقرات برتر است. کنترل فقط‌تصویری همان مدل از بازنمایی اصلی ضعیف‌تر است و حتی پس از افزودن عروق هم این خط لوله از مدل‌های اصلی جدا نمی‌شود. پس دو سطر جدول را باید دو آزمایش مستقل با یادگیرنده و رژیم آموزش متفاوت خواند و افزایش عددی بزرگ‌تر در سطر دوم شاهدی برای برتری آن معماری نیست (جدول شکل مربوط).

\begin{figure}[htbp]
  \centering
  \imgbox[1.0\textwidth]{figures/report/main/C3_vessel_complementarity_across_rgb_architectures.pdf}
  \caption{تفاضل سطح زیر منحنی از افزودن نمایش عروقی به دو بازنمایی تصویری متفاوت، با بازه‌ی اطمینان زوجی روی آزمون کانونی (1331 تصویر). این شکل برتری معماری دو‌ستون‌فقرات را نشان نمی‌دهد.}
  \label{fig:U21}
\end{figure}

\begin{figure}[htbp]
  \centering
  \imgbox[1.0\textwidth]{figures/report/appendix/A1a_attention_by_class.pdf}
  \caption{میانگین وزن توجه به شاخه‌ی \lr{ResNet} بر حسب کلاس واقعی در آزمون. این وزن‌ها فقط تشخیص توصیفی‌اند.}
  \label{fig:U22}
\end{figure}

\begin{figure}[htbp]
  \centering
  \imgbox[1.0\textwidth]{figures/report/appendix/A1b_attention_by_source.pdf}
  \caption{میانگین وزن توجه به شاخه‌ی \lr{ResNet} بر حسب منبع تصویربرداری. وابستگی شدید به منبع، تفسیر بالینی این وزن‌ها را نامعتبر می‌کند.}
  \label{fig:U23}
\end{figure}

\begin{figure}[htbp]
  \centering
  \imgbox[1.0\textwidth]{figures/report/appendix/A2a_attention_distribution_val.pdf}
  \caption{توزیع وزن توجه شاخه‌ی \lr{ResNet} روی اعتبارسنجی؛ دامنه‌ی صدک پنج تا نود و پنج با نشانگر میانه و میانگین.}
  \label{fig:U24}
\end{figure}

\begin{figure}[htbp]
  \centering
  \imgbox[1.0\textwidth]{figures/report/appendix/A2b_attention_distribution_test.pdf}
  \caption{توزیع وزن توجه شاخه‌ی \lr{ResNet} روی آزمون؛ دامنه‌ی صدک پنج تا نود و پنج با نشانگر میانه و میانگین.}
  
\end{figure}


پرسش دوم به وزن‌های توجه همین معماری دو‌ستون‌فقرات مربوط است. میانگین وزن ستون فقرات \lr{ResNet50} در اعتبارسنجی و آزمون به آستانه‌ی فروپاشی نزدیک نمی‌شود، ولی توزیع هر تصویر جداگانه تقریباً دوقطبی است و وزن‌های میانی کم دیده می‌شوند. یعنی مدل برای هر تصویر تقریباً به یکی از دو ستون فقرات متعهد می‌شود.

این وزن‌ها هم به کلاس و هم به منبع وابسته‌اند و وابستگی به منبع از ترتیب کلاس‌ها پیروی نمی‌کند. چون توزیع وزن این‌گونه به منبع تصویربرداری گره خورده است، نمی‌توان آن را توضیح بیماری‌محور خواند و به‌عنوان شاهد تفسیری ارائه کرد. این وزن‌ها در انتخاب مدل هم هیچ نقشی نداشتند؛ انتخاب مدل و توقف زودهنگام فقط با اعتبارسنجی انجام شد و این تشخیص ثانویه و اکتشافی است.




